% SOURCE_AUTHORITY: sga5_fr_workpass.tex, lines 4441--4494 (LF-normalized unit).
% SOURCE_SHA256: 791F4EFFC5E02832D5D77ED03518C8156D6F07E4C8238B03545DB93D883FBB28
A ideia da demonstração é que, nesse caso, $u$ e $v$ são imagens diretas de correspondências entre complexos concentrados nos pontos fechados e que a fórmula a estabelecer nada mais faz que exprimir a fórmula de Lefschetz para as projeções $x\to S$ e $y\to S$. Levando em conta 1.1.3, podem-se substituir os dados de 1.1 pelos seguintes: $X$ (resp. $Y$) é uma $S$-curva suave munida de um ponto fechado $x$ (resp. $y$), $A$ (resp. $B$) é uma $S$-curva separada, e $a:A\to X\times_S Y$, $b:B\to X\times_S Y$ são $S$-morfismos finitos tais que $A\times_{X\times_SY}B$ se reduza a um ponto fechado $z$ situado sobre $(x,y)$,
\[
u\in\Hom(a_1^*L,a_2^!M),\qquad v\in\Hom(b_2^*M,b_1^!L),
\]
com
\[
L=h_{1*}L',\quad M=h_{2*}M',\quad
L'\in\Ob D_{\mathrm{ctf}}(\{x\}),\quad M'\in\Ob D_{\mathrm{ctf}}(\{y\}),
\]
onde $h_1:\{x\}\to X$ e $h_2:\{y\}\to Y$ designam as inclusões canônicas. A questão é demonstrar que
\[
\langle u_z,v_z\rangle=\langle u,v\rangle_z,
\]
onde
\[
u_z\in\Hom(\R\Gamma(X,L),\R\Gamma(Y,M)),\qquad
v_z\in\Hom(\R\Gamma(Y,M),\R\Gamma(X,L))
\]
estão definidos como em 1.1, e $\langle u,v\rangle_z$ é o valor em $z$ do produto cup (III 4.2.7). Tem-se $\R\Gamma(X,L)=\R\Gamma_c(X,L)=L_x$, $\R\Gamma(Y,M)=\R\Gamma_c(Y,M)=M_y$, e $u_z$ (resp. $v_z$) nada mais é que a imagem direta $u_*$ (resp. $v_*$) de $u$ (resp. $v$) sobre $S$ no sentido de (III 3.7.6), como mostra o procedimento de cálculo (III 3.5 b)).

Por outro lado, se $h=h_1\times_Sh_2:(x,y)\to X\times_SY$ é a inclusão, e se $a':A'\to(x,y)$ e $b':B'\to(x,y)$ designarem as fibras de $a$ e $b$ em $(x,y)$, cada uma delas reduzida a um ponto, temos isomorfismos canônicos
\[
h_*a'_*a'{}^!(DL_x\lten_SM_y)
 \xrightarrow{\sim} a_*a^!(DL\lten_SM),
\]
\[
h_*b'_*b'{}^!(L_x\lten_SDM_y)
 \xrightarrow{\sim} b_*b^!(L\lten_SDM),
\]
e as flechas
\[
h_*:\Hom(a_1'{}^*L_x,a_2'{}^!M_y)\to \Hom(a_1^*L,a_2^!M),
\]
\[
h_*:\Hom(b_2'{}^*M_y,b_1'{}^!L_x)\to \Hom(b_2^*M,b_1^!L)
\]
(III 3.7.6) são isomorfismos. Portanto, as correspondências $u$ e $v$ se escrevem de modo único
\begin{equation}\tag{1}
u=h_*u',\qquad v=h_*v',
\end{equation}
com
\[
u'\in\Hom(a_1'{}^*L_x,a_2'{}^!M_y),\qquad
v'\in\Hom(b_2'{}^*M_y,b_1'{}^!L_x).
\]
Pela fórmula de Lefschetz (III 4.5.1) para $(h_1,h_2,A'\to A,B'\to B)$, tem-se
\begin{equation}\tag{2}
\langle u,v\rangle_z=\langle u',v'\rangle_z .
\end{equation}
Por outro lado, de acordo com (1), temos $u_*=u'_*$, $v_*=v'_*$, de onde, pela fórmula de Lefschetz para $\{x\}\to S$, $\{y\}\to S$,
\begin{equation}\tag{3}
\langle u_*,v_*\rangle=\langle u',v'\rangle_z .
\end{equation}
Levando em conta a observação anterior, a conclusão resulta da combinação de (2) e (3).
